\section{1. Resumen ejecutivo}

Enlace Broadcast Platform (EBP) se originó como respuesta a la necesidad de disponer de una solución propia para la distribución de la señal de EJTV. El desarrollo realizado ha permitido ampliar ese objetivo inicial hacia una arquitectura capaz de manejar varias señales, diferentes tecnologías de transporte y receptores externos, incorporando además una capa propia de supervisión operacional.

Actualmente EBP dispone de un servidor multimedia operativo, mecanismos de ingesta y distribución mediante SRT, RTMP, RTSP, HLS y WebRTC, y un Network Operations Center (NOC) que centraliza información sobre infraestructura, señales, sesiones, clientes, tráfico, eventos, alarmas e historial.

Las señales ENLACE y EJTV constituyen los casos principales de la plataforma. Cable Centro, ubicado en la Zona Sur de Costa Rica, representa el primer Reader o receptor real identificado dentro de la operación. IMPACT se utiliza como caso de prueba y diagnóstico técnico con una fuente real.

\begin{callout}
El avance alcanzado permite considerar a EBP no solamente como una solución puntual para EJTV, sino como una base técnica reutilizable para las necesidades de distribución audiovisual de Enlace.
\end{callout}
